% Appendix: Questions and Answers
% User annotations and interactive Q&A records

\section{阅读过程中的问题与解答}

% === 用户问答记录 ===
% 本节记录阅读过程中的问答内容

\section{什么是分布的核}\label{sec:qa:DistributionKernel}

\subsection{问题}
分布的核是什么概念？

\subsection{回答}
分布的\textbf{核（kernel）}是指概率密度/质量函数中\textbf{依赖于参数的核心部分}，不包括归一化常数。

\textbf{举例说明}：

Beta-Binomial 共轭中，后验分布：
\[
p(\theta|y) \propto \theta^{\alpha+y-1}(1-\theta)^{\beta+n-y-1}
\]

其中 $\theta^{\alpha+y-1}(1-\theta)^{\beta+n-y-1}$ 是\textbf{核}，而完整的 Beta 分布还需要归一化常数：
\[
p(\theta|y) = \frac{\Gamma(\alpha+\beta+n)}{\Gamma(\alpha+y)\Gamma(\beta+n-y)} \theta^{\alpha+y-1}(1-\theta)^{\beta+n-y-1}
\]

\textbf{为什么核重要}：

\begin{enumerate}
  \item \textbf{识别分布形式}：知道核是 $\theta^a(1-\theta)^b$ 形式，就能识别这是 Beta 分布
  \item \textbf{贝叶斯推断中省略 $p(y)$}：$p(y)$ 是归一化常数，不包含关于 $\theta$ 的信息，所以可以省略
  \item \textbf{简化推导}：在推导过程中只关注核，省去繁琐的归一化常数
\end{enumerate}

★ Insight ─────────────────────────────────────
这本质上是\textbf{信息局部性}思想：归一化常数的作用只是"让面积等于1"，它本身不携带关于参数的信息。真正"有意义"的概率质量都集中在核上。
─────────────────────────────────────────────────

% === QA 记录格式 ===
% 使用以下格式记录问答：
%
% \section{问题标题}
% \label{sec:qa:LabelKey}
%
% \subsection{问题}
% 用户的原始问题
%
% \subsection{回答}
% 回答内容...
%
% 如需详细推导，添加：
% \footnote{推导见附录 \cref{sec:xxx}}

\section{学习笔记}

\textit{在这里记录你的学习心得和思考。}

% === 学习笔记格式 ===
% \section{笔记标题}
% \label{sec:notes:LabelKey}
%
% 笔记内容...
